\centerline{\bf \Huge Третья МатДрака}




\problem{\textbf{1}} Найдите все решения ребуса: ИГРА = 108 $\times$ ИГ

\problem{\textbf{2}} Найдите все решения ребуса: 2$\times$МИР+3$\times$РИМ=2026

\problem{\textbf{3}} Найдите все решения ребуса: ГДЕ+КУДА=УДКА

\problem{\textbf{4}}{Разместите числа от 1 до 9, каждое ровно один раз,
в таблице $3\times3$ так, чтобы сумма чисел в каждом из четырёх квадратов
$2\times2$ делилась на 9.}

\Needspace{7\baselineskip}
\problem{\textbf{5}}{Разместите числа от 1 до 12, каждое ровно один раз,
в таблице $3\times4$ так, чтобы суммы во всех трёх строках были равны,
а суммы в четырёх столбцах были четырьмя последовательными натуральными
числами.}

\Needspace{7\baselineskip}
\problem{\textbf{6}}{Замостите клетчатую доску $6\times6$ доминошками
$1\times2$ (доминошки можно поворачивать) так, чтобы в каждой строке лежала
ровно одна горизонтальная доминошка, а в каждом столбце~--- ровно две
вертикальные.}

\Needspace{7\baselineskip}
\problem{\textbf{7}}{В четырёхзначном числе $\overline{3a5b}$ цифры $a$ и
$b$ поменяли местами и получили число $\overline{3b5a}$. Известно, что
$\overline{3a5b}$ делится на 24(на 3 и на 8), а $\overline{3b5a}$ делится на 18(на 2 и на 9).
Найдите цифры $a$ и $b$.}

\Needspace{6\baselineskip}
\problem{\textbf{8}}{Расположите цифры $0,1,4,8$, использовав каждую ровно
один раз, так, чтобы получилось четырёхзначное число, делящееся на 88(на 8 и на 11).
Найдите все такие числа.}

\Needspace{6\baselineskip}
\problem{\textbf{9}}{Сколько пятизначных чисел можно составить из цифр
$0,2,3,4,6$, использовав каждую ровно один раз, если число должно делиться
на 6(на 2 и на 3)?}


\problem{\textbf{10}}{На фестивале мастерских Лида, Миша, Оля и Тимур
посетили по одному разному мастер-классу: гончарное дело, печать,
робототехника и витраж. Каждый сделал один сувенир, причём все сувениры
были разными: чашка, открытка, фонарик и сова. Известно, что:
\begin{enumerate}[label=\arabic*),leftmargin=8mm,itemsep=0pt,topsep=1mm]
    \item Лида не делала чашку;
    \item Тимур сделал сову;
    \item Тимур не был на гончарном мастер-классе;
    \item Оля занималась печатью;
    \item на мастер-классе по робототехнике сделали фонарик;
    \item на мастер-классе по печати изготовили открытку.
\end{enumerate}
Определите, на каком мастер-классе был каждый из ребят и какой сувенир он
сделал.}

\Needspace{18\baselineskip}
\problem{\textbf{11}}{Перед походом пять ребят~--- Аня, Боря, Вера, Глеб и
Дина~--- собрали пять коробок. Коробки поставили в ряд и пронумеровали слева
направо числами от 1 до 5. Каждый ребёнок собрал ровно одну коробку. В
коробках лежали разные предметы: бинокль, компас, лупа, фонарь и фляга.
Известно, что:
\begin{enumerate}[label=\arabic*),leftmargin=8mm,itemsep=0pt,topsep=1mm]
    \item фляга лежит в коробке №~4;
    \item Динина коробка стоит левее коробки с биноклем;
    \item коробка с лупой стоит левее Дининой коробки;
    \item Глебова коробка стоит непосредственно слева от коробки с компасом;
    \item Глебова коробка стоит непосредственно слева от Аниной коробки;
    \item Боря положил в свою коробку флягу.
\end{enumerate}
Определите, кто собрал каждую коробку и какой предмет в ней лежит.}


\Needspace{8\baselineskip}
\problem{\textbf{12}}{На фестивале 36 школьников посетили хотя бы одну из
трёх площадок: «Шифры», «Мосты» и «Зеркала». За каждую посещённую площадку
школьник получал ровно одну отметку; всего выдали 62 отметки. Ровно одну
площадку посетили 15 школьников. Сколько школьников посетили все три
площадки?}

\Needspace{10\baselineskip}
\problem{\textbf{13}}{В отряде 42 школьника. Шахматами занимаются 20
человек, моделированием~--- 24, программированием~--- 15. Каждый программист
занимается также моделированием. Шахматами и моделированием занимаются 9
человек, а шахматами и программированием~--- 6 человек. Сколько школьников
занимается только моделированием (не шахматами и не программированием)?
Сколько не занимается ни одним из этих трёх дел?}


\problem{\textbf{14}} Сколько существует натуральных чисел, меньших тысячи, которые не делятся ни на 5, ни на 7?

\problem{\textbf{15}} Туристическая фирма провела акцию: «Купи путёвку в Египет,
приведи четырёх друзей, которые также купят путёвку, и получи стоимость путёвки обратно».
За время действия акции 13 покупателей пришли сами, остальных привели друзья. Некоторые
из них привели ровно по 4 новых клиента, а остальные 100 не привели никого. Сколько туристов
отправились в Страну Пирамид бесплатно?
